\documentclass[11pt]{article}

\usepackage[T1]{fontenc}
\usepackage[utf8]{inputenc}
\usepackage[margin=1in]{geometry}
\usepackage{amsmath, amssymb}
\usepackage{textcomp}
\usepackage{setspace}
\usepackage[hidelinks]{hyperref}
\usepackage{url}

\newcommand{\MISSING}{\textbf{[MATH MISSING]}}

\setlength{\parindent}{0pt}
\setlength{\parskip}{0.8em}
\onehalfspacing

\title{FIREFate Methods}
\date{}

\begin{document}

\section*{Methods}

\subsection*{FIREFate overview}

The FIREFate workflow consists of: (1) inferring cellular programs underlying cell
states using an interpretable machine learning model along with characterizing their
phenotypic roles through gene set enrichment analysis supported by large language
models; (2) constructing state-specific static GRNs using sc-multiomics data or
scRNA-seq data alone; (3) defining highly specific TF-centric regulons by integrating
the inferred cellular programs with the static GRN and calculating TF-centric regulon
enrichment; (4) constructing state-specific dynamic GRNs using sc-multiomics data;
(5) identifying state-specific TF expression and regulatory dynamics, including
episodic enrichment patterns; (6) identifying state-predisposing cellular programs
that bias cell-fate decisions using an interpretable machine learning model applied to
TF-KO versus unperturbed control cell populations; and (7) determining TFs enriched
within these predisposing programs through integration of static and dynamic GRNs.
We implemented and tested FIREFate in Python and R (versions xx and xx) and designed
% TODO: fill in Python / R version numbers.
it for use in Jupyter Notebook and RStudio environments. FIREFate code is open source
and available on GitHub (\url{https://github.com/sachha-naksha/FIREFate.git}), along
with detailed descriptions of functions and tutorials. Additionally, we provide a
user-friendly web application (\url{https://pitt-csi.shinyapps.io/firefate/}) that
facilitates FIREFate analyses and enables interactive exploration of the results
presented in this manuscript.

\subsection*{Inferring cellular programs using interpretable machine learning}

In the first step, the cellular programs underlying cell states using the scRNA-seq
data are identified. To achieve this, we employ an interpretable machine learning
method called SLIDE to identify the significant cellular programs (CPs) underlying
outcomes of interest from high-dimensional omic datasets [ref to SLIDE]. This method
is beneficial as it discovers necessary and sufficient CPs to distinguish the states.
SLIDE is a regression-based model that identifies significant CPs ($Z$) capturing
linear and nonlinear relationships between observed data ($X$) and the phenotype of
interest ($Y$: in our case binary values regarding each state of interest):

\begin{equation}
X' = AZ + E
\end{equation}

$X_{n \times p}$ is scRNA-seq expression data with $n$ cells and $p$ genes.
$X_{n \times p}$ is decomposed into two factors, $A_{p \times K}$ and $Z_{K \times n}$,
with an error term $E$. $A$ is the allocation matrix and represents the membership of
each feature to a cellular program. $Z$ is the CP matrix and represents a
lower-dimensional representation ($K < p$) of the input data in latent space. Then,
using a regression model applied on CPs, the linear part of the model is generated,
which represents the significant standalone CPs:

\begin{equation}
LP = \sum_{j \in S_{1}} \beta_{j} Z_{j} + \epsilon_{1}
\end{equation}

Here $LP$ is the linear part of the model, and $S_{1}$ is the set of significant
standalone CPs. Then, incorporating the nonlinear relationships, the method
identifies significant interacting CPs:

\begin{equation}
NP_{j} = \beta_{j} Z_{j} + \sum_{i} C_{ij}\, Z_{i} \odot Z_{j} + \epsilon_{2}
\qquad \text{where } i \neq j,\; j \in S_{1},\; i \in \{1 \dots K\},\;
\text{and } Z_{i} \odot Z_{j} \in S_{2}
\end{equation}

Here $\beta \in \mathbb{R}^{1 \times K}$ and $C_{ij}$ is the effect size of the
interaction term between the two CPs $Z_{i}$ and $Z_{j}$. $S_{2}$ is the set of putative
interactions involving the standalone significant CPs. These interactions arise from the
model's non-linear architecture, capturing combinatorial regulation and helping improve
interpretability and predictive performance.

Before applying SLIDE, the single-cell RNA sequencing (scRNA-seq) data were
preprocessed using the filtering and preprocessing functions provided in the SLIDE
library (\url{https://github.com/jishnu-lab/SLIDE}) and FIREFate GitHub repository
(\url{https://github.com/jishnu-lab/xxxx}), following standard best practices. First,
% TODO: replace jishnu-lab/xxxx with the final FIREFate repo URL (currently
% https://github.com/sachha-naksha/FIREFate) -- these two URLs disagree in the draft.
genes with zero unique molecular identifier (UMI) counts across all cells were removed,
and mitochondrial and ribosomal genes were excluded from further analysis. Second,
sparsity filtering was performed using the \texttt{zeroFiltering} function in the SLIDE
library to remove genes with more than \texttt{col\_thresh} zero values across cells.
To identify the cellular programs (CPs), SLIDE was applied to two distinct single-cell
RNA sequencing (scRNA-seq) datasets to extract discriminatory regulons between
predefined cell states in B cells and T cells. SLIDE requires two input cell populations
to learn latent regulatory programs that differentiate one state from another. For the
B cell system, the method was applied to PB and GC B cells using both scMultiome data
(Day 6) and scRNA-seq data. For the T cell system, the model was applied to Tex-term
and Tex-KLR cells. A summary of the training conditions---including experimental groups,
collection time points, and model parameters---is provided in Supplementary Table 1.
% TODO (highlighted in draft): confirm Supplementary Table 1 exists / renumber.

We also used this approach to get the predisposing cellular programs from single-cell
Perturb-seq data collected on Day 6 following transcription factor (TF) KOs. The
training dataset for B cell included KO conditions for key TFs implicated in B cell
differentiation (PRDM1, BATF, SPIB, IRF4, and IRF8) alongside unperturbed control cells.
For T cell system we applied the same approach on single-cell Perturb-seq data of Ikzf1
and Ets1 KO cells against unperturbed control cells. A summary of the training
conditions---including experimental groups, and model parameters---is provided in
Supplementary Table 2.
% TODO (highlighted in draft): confirm Supplementary Table 2 exists / renumber.

\subsection*{Interpreting the cellular programs using LLM-based gene set enrichment analysis}

To derive biologically meaningful modules and corresponding labels for each CP, we
employed a large language model (LLM)-based gene set enrichment approach on the
top-ranked genes within each CP. GSAI grouped genes into functional programs and
generated module annotations. Using the prompt provided in Supplementary File S1, we
% TODO (highlighted in draft): attach Supplementary File S1 (GSAI prompt).
supplied the sets of top-ranked genes and asked ChatGPT 5.1 to annotate them with the
most relevant biological processes.

\subsection*{State-specific static GRN inference}

GRNs were constructed separately for B cell and T cell datasets, each containing
distinct sub-cell types. For the B cell dataset, GRNs were obtained for the PB and GC
cells from [ref Nick's]. For the T cell dataset, GRNs were inferred for the Tex-term
and Tex-KLR cells using the CellOracle framework, with its murine ATAC-seq inferred GRN
serving as the base GRN. For each cell type (B or T), GRNs corresponding to the two
relevant cell-subtypes were first subset to top 10{,}000 edges. These subtype-specific
GRNs were then merged into a fusion GRN that represents the combined regulatory
landscape across sub-cell types within that dataset. During the merging process, if an
edge (TF--target interaction) was present in both subtype GRNs, the edge with the higher
weight (indicative of stronger regulatory potential) was retained. If an edge was unique
to one subtype, it was included as-is in the fusion GRN. Following GRN fusion, each edge
was categorized based on its weight relative to the distribution of edge weights in the
respective fusion network: strong regulatory edges, edge weight in the top 10th
percentile ($\geq$ 90th percentile), and weak regulatory edges, edge weight below the
90th percentile. This strategy enabled comprehensive subtype-integrated GRNs across
functionally distinct immune subpopulations.

\subsection*{Inferring regulons by combining cellular programs and GRNs}

To identify TFs with regulatory programs enriched in cell subtype discrimination
specific regulons, we performed enrichment analyses using previously constructed fusion
GRNs alongside regulons inferred from SLIDE. By overlapping GRN topology with both
standalone and interacting regulons, we assembled a comprehensive set of TF--target gene
relationships. These interactions were then evaluated through two levels of enrichment
analyses. The order-1 (single-TF) enrichment analysis, in which for each regulon, TFs
either directly included in the regulon or connected as first-degree neighbors (based on
GRN topology) were considered. To assess their regulatory importance, we quantified how
many genes each TF regulates strongly within the regulons compared to its global
regulatory profile in the fusion GRNs. For each TF, the enrichment score was computed as:

\begin{equation}
\text{Enrichment Score} \;=\;
\frac{
  \dfrac{\text{Number of \textbf{strong} TF} \rightarrow \text{gene edges within the regulon}}
        {\text{Total number of genes in the regulon}}
}{
  \dfrac{\text{Number of \textbf{strong} TF} \rightarrow \text{gene edges in the fusion GRNs}}
        {\text{Total number of genes in the GRN}}
}
\end{equation}

The enrichment score thus quantifies the TF's overrepresented strong regulatory influence
for given regulons. A hypergeometric test was used to determine statistical significance,
and transcription factors with significantly enriched regulons were identified as
candidate key regulators.

The second set of regulons are the order-2 (TF-pair) enrichment analysis. To capture
potential combinatorial regulation, we extended the analysis to transcription factor
pairs (TF pairs). Like order-1, we evaluated TF pairs that were either members of a
regulon or connected to regulon genes as first-degree neighbors from fusion GRNs. For
each TF pair, we identified the set of common downstream genes that were co-regulated
within the regulon. Only interactions where at least one of the TF-to-gene edges were
classified as strong were considered; weak--weak pairs were excluded. The observed number
of co-regulated genes within the regulon was then compared to the expected overlap based
on the global GRN.

\begin{equation}
\text{Enrichment Score}_{\text{order-1}} \;=\;
\frac{E^{\text{strong}}_{\text{reg}} / G_{\text{reg}}}
     {E^{\text{strong}}_{\text{GRN}} / G_{\text{GRN}}}
\end{equation}

\begin{equation}
\text{Enrichment Score}_{\text{order-2}} \;=\;
\frac{P_{\text{reg}} / G_{\text{reg}}}
     {P_{\text{GRN}} / G_{\text{GRN}}}
\end{equation}

where $E^{\text{strong}}_{\text{reg}}$ and $E^{\text{strong}}_{\text{GRN}}$ are the
numbers of strong TF$\rightarrow$gene edges within the regulon and within the fusion
GRN respectively; $P_{\text{reg}}$ and $P_{\text{GRN}}$ are the numbers of genes
co-regulated by a TF pair with at least one strong edge (i.e. (strong, strong),
(strong, weak) or (weak, strong); weak--weak pairs excluded) within the regulon and
within the fusion GRN respectively; and $G_{\text{reg}}$ and $G_{\text{GRN}}$ are the
total numbers of genes in the regulon and in the GRN.

To reduce spurious associations, we applied two filtering criteria:
(1) TF pairs must co-regulate more than one gene within the regulon, and (2) TF pairs
must co-regulate at least four genes across the full GRN. Statistical enrichment was
again assessed using a hypergeometric test, and TF pairs with significant enrichment
($p < 0.05$) that passed both filters were retained as candidate combinatorial
regulators contributing to subtype-specific transcriptional programs.

\subsection*{Benchmarking of FIREFate cellular programs}

To enable direct comparison of order-1 transcriptional programs enriched by FIREFate,
TF-centric regulons, we employed SCENIC+ to generate state-specific regulons
distinguishing GC and PB states in the B-cell system (Supplementary Fig.~1), and Tex-term
and Tex-KLR states in the T-cell system (Supplementary Fig.~2). Our workflow proceeded as
follows: SCENIC+ was first applied to single-cell multi-omic data of both systems to
infer regulons differentially active between the specified cell states. Regulon activity
and specificity were quantified using the \textit{eRegulon} score. To ensure comparability
across analyses, we retained only those SCENIC+-derived regulons whose TFs overlapped with
significantly enriched order-1 TFs identified through our regulon and fusion GRN
analyses. To compare the specificity and sensitivity of those TF-centric regulons, we
compiled the target genes of each transcription factor within the selected SCENIC+
regulons, as well as those defined by the regulon and fusion GRN frameworks. We then
calculated the absolute fold changes in gene expression between the two cell state
subtypes to assess the differential regulatory impact across methodologies.

\subsection*{Imputation of multi-ome dataset}

PALANTIR was used to impute the sc-RNA counts of the B-cell multi-omic (scRNA+ATAC)
dataset. Following standard preprocessing, highly variable features were selected for
PCA dimensionality reduction. 10 PC components were then used to get the Diffusion Maps
eigenvector basis and cellular neighborhood for the B-cell states. Palantir's built-in
MAGIC algorithm was used to impute the gene expression of the highly variable genes
($n = 3018$). The imputed dataset was then filtered to only contain cell states of
interest for trajectory inference, dropping `Naive', `8\_UD\_2', and `11\_UD\_3' states.

\subsection*{Inferring pseudo-time trajectory}

To reconstruct single-cell differentiation trajectory from our dataset of 28,494 cells
in the B-cell system, we applied STREAM on the MAGIC imputed gene-expression layer with
`all' highly variable genes. A low-dimensional manifold of cellular neighborhood
reflective of the continuous differentiation process was learnt, and branching
trajectories were inferred using STREAM's standard locally linear embeddings (SE) and
elastic principal graph framework, respectively. We applied STREAM on 2088 cells in the
T-cell system from the non-targeted control population of the RBPJ-KO Perturb-seq
dataset, without imputing gene-expression. Raw gene-expression counts were used to filter
cells and features from the T-cell dataset, which were then scaled and normalized. 370
variable genes were selected for dimensionality reduction, by fitting a loess curve on
the standard deviation and mean values of features with a fraction of 0.01 and 80th
percentile. The parameters for fitting the elastic graph on both the T-cell and B-cell
trajectories were: \texttt{epg\_alpha} = 0.02, \texttt{epg\_mu} = 0.07,
\texttt{epg\_lambda} = 0.02, \texttt{epg\_trimmingradius} = 2. The method generated
bifurcating and linear trajectories, capturing the underlying transient cellular states
in both the differentiating B cells and state-switching T cells, and assigned pseudotime
values to each cell to order them on the temporal scale. We used STREAM's visualization
modules to inspect trajectory branches, ordering of cell states based on their sequential
transitions, and gene expression changes along pseudotime branches.

\subsection*{Defining dynamic TF force on a target gene}

% TODO: section not yet written.
% Placeholder note from the draft: "SDE adaptation to include atac+rna / only-rna with
% base GRN + equation derived empirically."
%
% Suggested scaffold matching the implementation (signed force = sign of the fitted
% coefficient x its magnitude x the TF's expression at that point in pseudotime).
% Uncomment and adapt once the SDE formulation is finalised:
%
% For a transcription factor $f$ acting on target gene $g$ at pseudotime $t$, we define
% the dynamic TF force as
% \begin{equation}
% F_{f \rightarrow g}(t) \;=\; \operatorname{sign}\!\big(\beta_{f \rightarrow g}(t)\big)
% \cdot \big|\beta_{f \rightarrow g}(t)\big| \cdot x_{f}(t),
% \end{equation}
% where $\beta_{f \rightarrow g}(t)$ is the regulatory coefficient of the edge
% $f \rightarrow g$ in the episodic GRN local to pseudotime $t$, and $x_{f}(t)$ is the
% (smoothed) expression of $f$ at $t$. Positive and negative values of
% $F_{f \rightarrow g}(t)$ correspond to activating and repressive force respectively.

SDE adaptation to include atac+rna / only-rna with base GRN + equation derived
empirically.

\subsection*{Construction of gene regulatory networks across episodes of state transition trajectories}

% TODO: section not yet written.

\subsection*{Enriching for dynamic TF activity in episodic GRNs}

% TODO: section incomplete -- sentence below is truncated in the draft.
The trajectory branches for both the datasets were broken into episodes of invariant TF
activity. T cells have a linear

\subsection*{Transient regulatory activity analysis}

% TODO: section not yet written.

\subsection*{Validation in-vivo dataset??}

% TODO: section not yet written.

\subsection*{Chromatin binding inference across pseudotime trajectories}

% TODO: section not yet written.

\subsection*{Deriving phases of gene-regulation governing state switching of cells}

% TODO: section not yet written.

\subsection*{Utilization of optimal transport for validating dynamic TF activity}

% TODO: section not yet written.

\subsection*{Fate prediction of early activated B cells}

To predict cell fate bias toward the final lineages, we first evaluated whether
transcriptomic features that distinguish TF-KO cells from controls could also serve as
informative predictors. We identified key transcriptomic features that differentiate
PRDM1- and IRF4-perturbed cells from wild-type controls by training a Lasso logistic
regression classification model (\texttt{LogisticRegression(penalty='l1',
solver='liblinear', max\_iter=10000, random\_state=42)}). The trained model was then
applied to early ABCs from the multi-omic dataset (days 2 and 4) using the
\texttt{lasso.predict(...)} function to generate predicted probabilities for each cell.
Because the response variable is binary, we assigned predicted fates to early ABCs based
on whether their predicted probability was above or below a threshold of 0.5.

The same procedure was repeated using only the differentially expressed genes from the
training datasets. We first performed differential expression analysis between KO and
control groups and retained genes with an absolute $\log_{2}$ fold change greater than 1
and an adjusted $p$-value below 0.05. The single-cell gene expression matrix was then
subset to include only these genes, and a Lasso logistic regression model was trained on
this reduced feature set. The trained model was subsequently applied to early ABCs from
the multi-omic dataset (days 2 and 4) using the same prediction function to generate and
project fate probabilities.

We then used the inferred cellular programs from single-cell TF Perturb-seq data in the
B cell system. FIREFate projected these programs onto ABCs (day 2 and day 4) from
single-cell multiome data to estimate their transcriptional predisposition toward one of
two downstream fates (GC or PB). The latent cellular programs defined a score for each
cell, representing the degree to which its expression profile aligned with either the KO
or control state. Using these scores, we assessed the discriminative power of each model
by examining how well the score distributions separated KO and control cell populations.
For each model, we defined an optimal threshold based on the distribution that most
effectively distinguished these two groups. We then applied the same trained models to
the ABCs at day 2 and day 4 multiome data using the \texttt{predZ} function in the SLIDE
package, generating predicted scores for uncommitted cells. Using the thresholds defined
from the Perturb-seq models, FIREFate assigned a predicted fate score to each ABC. Cells
with predicted scores above or below the KO-control threshold were classified as
transcriptionally predisposed toward different downstream fates.

To evaluate the quality of these predictions, we trained a Lasso logistic regression
classifier to assess whether the predicted fate could be distinguished from the true fate
labels (GC and PB cells from the single-cell multiome dataset). Two scenarios were
defined: first, if a predicted fate could not be distinguished from the true identity
based on transcriptional features, the cell was assigned to the same identity.
Conversely, if the predicted fate could be confidently discriminated from the true
identity, we inferred that \textbf{FIREFate} did not specify the cell as belonging to the
predicted fate. This framework produced classification loss scores that quantified the
predictive confidence of \textbf{rollback analysis}, thereby enabling evaluation of the
model's performance in capturing early fate bias.

\subsection*{Gene set enrichment analysis using ESCAPE}

To interrogate transcriptional heterogeneity and quantify the enrichment of gene
expression signatures at the single-cell level, we employed the ESCAPE (Enrichment of
Single Cell Analysis by Pathway Expression) framework (Borcherding et al., Bioconductor,
\url{https://bioconductor.org/packages/escape}). ESCAPE is a single-cell gene set
enrichment analysis (GSEA) tool specifically designed to leverage the high dimensionality
and variability inherent to scRNA-seq data. By integrating predefined gene sets or custom
marker lists, ESCAPE calculates enrichment scores for individual cells, enabling
sensitive detection of transcriptional programs and cellular states within heterogeneous
populations. In our study, we applied ESCAPE to assess whether the transcriptional
profiles of predicted ``KO-like'' subsets reflected expected lineage biases. Enrichment
was calculated using canonical marker genes representative of progenitor-like exhausted
T cells (Tpex: \textit{Tcf7}, \textit{Slamf6}, \textit{Myb}, \textit{Sell},
\textit{Bach2}) and terminally exhausted T cells (Tex: \textit{Havcr2}, \textit{Pdcd1},
\textit{Entpd1}, \textit{Cd38}, \textit{Cd244a}, \textit{Cxcr6}, \textit{Ifng}). These
analyses allowed us to systematically compare enrichment patterns across subsets and gain
insights into the differentiation status and functional states of the populations under
investigation.

\end{document}